% GDZ PPN236018647,native589lower–590upper;HKUST590–591.
\NumberedItem{118}
\noindent\hspace*{\parindent}Methodus quinta theorema fundamentale demonstrandi se obtulit adiumento theorematis elegantissimi theoriae sectionis circuli, puta
\clearpage\SourceLeaf{590}\ArithmeticHead{561}{R}{1801 APR. 6—1801 SEPT. MEDIO.}
\[
\begin{array}{r@{\;}cccc@{\;}l}
\displaystyle\sum\left.\begin{matrix}\sin\\\cos\end{matrix}\right\}\frac{nn}{a}P=
&\multicolumn{1}{c|}{\begin{gathered}+\surd a\\+\surd a\end{gathered}}
&\multicolumn{1}{c|}{\begin{gathered}0\\+\surd a\end{gathered}}
&\multicolumn{1}{c|}{\begin{gathered}0\\0\end{gathered}}
&\begin{gathered}+\surd a\\0\end{gathered}&\\[5pt]
\text{prout }a\equiv&0&1&2&3&(\mathrm{mod.}\ 4)
\end{array}
\]
\noindent substituendo pro \(n\) omnes numeros a 0 usque ad \((a-1)\).
\par\noindent\hfill Bruns[vigae, 1801] Mai. medio.
\par\medskip\begingroup\small
\noindent\hspace*{\parindent}Vergl. hierzu die oben S. 23 abgedruckte Aufzeichnung \textit{Theorema novissimum pulcherrimum}. Siehe auch die Notiz Nr. 123 vom 30. August 1805.
\par\noindent\hfill\textsc{Klein.}
\par\endgroup
